\section{Introduction}
\label{sec:intro}

Aerial scene classification assigns a land-use or land-cover label, such as \emph{Port}, \emph{Railway} or \emph{River}, to an overhead image, and it underpins many mapping and monitoring tasks.
Two things make it hard.
Classes that share colour and texture must be told apart: on the SkyView data used here, a hand-crafted pipeline most often mistakes \emph{Highway} for \emph{Railway} and \emph{Lake} for \emph{River}.
And real collections are rarely balanced, so some classes come with far fewer training images than others~\cite{he2009imbalanced}.

The usual comparison pits two families of methods against each other.
Hand-crafted pipelines summarise an image with colour, texture or keypoint statistics and train a shallow classifier on top~\cite{swain1991color,ojala2002lbp,lowe2004sift,yang2010bovw}, while ImageNet-pretrained CNNs are fine-tuned end to end~\cite{nogueira2017towards,pires2020transfer}.
Ensembles are often suggested as a way to get the best of both~\cite{dietterich2000ensemble,sagi2018ensemble}, and augmentation or loss re-weighting is the standard response to imbalance~\cite{buda2018imbalance,zhang2023longtail}.
These comparisons are easy to get wrong, though.
If each method is tuned on its own split, or hyperparameters are chosen on the test images, a results table no longer separates the effect of the method from the effect of the protocol.

We ran into exactly this problem.
An earlier version of this study, a course project, tuned on the test set, mixed numbers from different splits, and broke voting ties by member order.
That last choice is subtle: with two voters every disagreement is a tie, so a vote that favours the first-listed member simply returns that member's predictions.
The ensembles looked as if they collapsed to their weakest member when in fact they were never combining anything.

This paper repeats the comparison under a protocol designed to avoid these traps.
We hold out one balanced test set and one balanced validation set, and derive three training sets from the remaining images that differ only in class distribution: balanced, long-tailed with an imbalance ratio of 16, and rebalanced by augmenting the minority classes (\cref{sec:protocol}).
Every hyperparameter and stopping epoch is chosen on the validation set.
On each training set we train four hand-crafted pipelines and two CNNs, \rn and \eb, the latter with a linear-probe-then-fine-tune schedule (\cref{sec:methods}), and we combine them by soft voting and by hard voting with a probability-based tie-break.
All comparisons are made on the same 2{,}400 test images, with bootstrap confidence intervals and paired McNemar tests.

Our contributions are:
\begin{itemize}
  \item a controlled protocol and re-implementation for comparing hand-crafted features, fine-tuned CNNs, ensembles and rebalancing strategies on SkyView (\cref{sec:protocol,sec:methods});
  \item measurements of how far apart the two families are on balanced data and how much each loses under a 16:1 long tail, broken down by head, middle and tail classes (\cref{sec:res_base,sec:res_imbalance});
  \item evidence that averaging the two CNNs significantly improves on either network in all three settings while adding hand-crafted members does not, together with a diagnosis of the tie-breaking pitfall (\cref{sec:ensembles,sec:res_ensembles});
  \item a comparison of minority-class augmentation with inverse-frequency loss weighting, which shows that augmentation helps the hand-crafted pipelines but not the CNNs (\cref{sec:res_imbalance}).
\end{itemize}
